\section{The finite fillings and the compact complex threefold}

\subsection{Free logarithmic transforms}
For $j=1,2$, choose a $g_j$-invariant disc $\Delta_j\ni z_j$ and a
coordinate $s_j$ such that
\[
 s_j(z_j)=0,
 \qquad
 s_j\circ g_j=e^{-2\pi i/m_j}s_j.
\]
The quotient coordinate $t_j$ on $D_j=\Delta_j/\langle g_j\rangle$ is
defined by $t_j\circ\pi=s_j^{m_j}$.  Let
\[
 J'_j=(\Delta_j\times\C^2)/\Lambda
\]
be the restriction of the marked torus family.  Define
\begin{equation}\label{eq:log-action}
 \wt g_j^{\mathrm{log}}(z,\zeta)
 =\left(g_jz,R_j(z)\zeta+\frac1{m_j}\Pi(g_jz)v_j\right).
\end{equation}
The identity
$R_j(z)\Pi(z)=\Pi(g_jz)A_j$ implies
\[
 \wt g_j^{\mathrm{log}}\,t_\lambda
 =t_{A_j\lambda}\,\wt g_j^{\mathrm{log}},
\]
so \eqref{eq:log-action} descends to $J'_j$.  In the smooth
trivialization
\[
 \Delta_j\times(\Lambda\otimes\R)/\Lambda\longrightarrow J'_j,
 \qquad (z,x)\longmapsto[(z,\Pi(z)x)],
\]
it is
\begin{equation}\label{eq:log-action-flat}
 (z,x)\longmapsto(g_jz,A_jx+v_j/m_j).
\end{equation}

\begin{lemma}[Fixed points of an affine torus map]\label{lem:affine-fixed}
Let $A\in\GL(\Lambda)$ have finite order and let $b\in\Lambda\otimes\Q$.
The affine map $x\mapsto Ax+b$ of $(\Lambda\otimes\R)/\Lambda$ has a fixed
point if and only if
\[
 \phi(b)\in\Z
 \quad\text{for every }\phi\in V\text{ satisfying }\phi\circ A=\phi.
\]
\end{lemma}

\begin{proof}
A fixed point exists exactly when
$b\in(A-I)(\Lambda\otimes\R)+\Lambda$.  Necessity is immediate.  For the
converse, average the powers of $A$ to obtain the rational projection
$p$ onto $\ker(A-I)_\R$ with kernel $\im(A-I)_\R$.  The lattice
$p(\Lambda)$ is dual to the invariant lattice
$\{\phi\in V:\phi A=\phi\}$.  The integrality assumption therefore gives
$pb=p\lambda$ for some $\lambda\in\Lambda$, so
$b-\lambda\in\im(A-I)_\R$.
\end{proof}

\begin{theorem}[The two logarithmic fillings]\label{thm:finite-fillings}
The map \eqref{eq:log-action} generates a free cyclic action of order
$m_j$ on $J'_j$.  Its quotient
\[
 N_j=J'_j/\langle\wt g_j^{\mathrm{log}}\rangle
\]
is a smooth complex threefold, and $s_j^{m_j}$ descends to a proper
holomorphic map
\[
 f_j:N_j\longrightarrow D_j
\]
with divisor
\[
 f_j^*(p_j)=m_jS_j,
\]
where $S_j$ is a smooth compact surface.

Over the punctured disc there is a biholomorphism
\[
 G_j:N_j\setminus S_j\longrightarrow J|_{D_j^*}.
\]
Under this biholomorphism the zero section of $J$ is represented on the
cover by
\begin{equation}\label{eq:zero-section-log}
 \varsigma_j(z)=-\frac{\log s_j(z)}{2\pi i}\Pi(z)v_j.
\end{equation}
Following this section around $s_j$ once counterclockwise changes its
lift by $-\Pi(z)v_j$.
\end{theorem}

\begin{proof}
Because $A_jv_j=v_j$, the $m_j$-th power of
\eqref{eq:log-action-flat} is translation by $v_j\in\Lambda$, hence the
identity.  A nontrivial power can have a fixed point only over $z_j$.
A direct calculation gives
\[
 V^{T_1}=V^{T_1^2}=\langle\gamma,\psi_1\rangle,
 \qquad \psi_1=2u+w+3\delta,
\]
and
\[
 V^{T_2}=V^{T_2^2}=V^{T_2^3}=\langle\gamma,\psi_2\rangle,
 \qquad \psi_2=u+w+2\delta.
\]
The $k$-th power of \eqref{eq:log-action-flat} has translation part
$kv_j/m_j$.  For $v_1=\eps$ one has
$\gamma(v_1)=1$ and $\psi_1(v_1)=0$; hence for $k=1,2$ the value
$k\gamma(v_1)/3$ is not integral.  For $v_2=-\eps'$ one has
$\gamma(v_2)=-1$ and $\psi_2(v_2)=0$; hence for $k=1,2,3$ the value
$-k/4$ is not integral.  The fixed-point criterion in
\cref{lem:affine-fixed} therefore excludes fixed points for every
nontrivial power.

A free finite holomorphic action has a smooth quotient.  The function
$s_j^{m_j}$ is invariant, properness follows from properness of
$J'_j\to\Delta_j$ and finiteness of the quotient, and its divisor is
$m_j$ times the reduced central fibre.

On $\Delta_j^*$ define the section
\[
 \sigma_j(z)=\frac{\log s_j(z)}{2\pi i}\Pi(z)v_j.
\]
Changing the logarithm changes $\sigma_j$ by a period, so it is well
defined on the torus family.  Translation by $\sigma_j$ conjugates
$\wt g_j^{\mathrm{log}}$ to the linear deck transformation $\wt g_j$.
Indeed
$s_j(g_jz)=e^{-2\pi i/m_j}s_j(z)$ and
$R_j(z)\Pi(z)v_j=\Pi(g_jz)A_jv_j=\Pi(g_jz)v_j$, so the translation terms
cancel.  This gives $G_j$.  The zero section in $N_j$-coordinates is
$-\sigma_j$, proving \eqref{eq:zero-section-log} and the stated drift.
\end{proof}

\begin{proposition}[Topology of the logarithmic pieces]\label{prop:pi1-Nj}
The radial contraction of the $s_j$-disc is equivariant and induces a
deformation retraction $N_j\simeq S_j$.  Its fundamental group has the
presentation
\[
 \pi_1(N_j)=
 \left\langle \Lambda,\wh g_j\ \middle|\
 \wh g_j\lambda\wh g_j^{-1}=A_j\lambda,
 \ \wh g_j^{m_j}=t_{v_j}
 \right\rangle.
\]
For the boundary mapping torus $M_j=f_j^{-1}(|t_j|=\rho)$ one has
$\pi_1(M_j)=\Lambda\rtimes_{A_j}\Z\langle\wh g_j\rangle$, and the
kernel of $\pi_1(M_j)\to\pi_1(N_j)$ is normally generated by
$\wh g_j^{m_j}t_{-v_j}$.
\end{proposition}

\begin{proof}
Before quotienting, the local model is smoothly
$A_0^{(j)}\times\Delta_{s_j}$, where
$A_0^{(j)}=\C^2/L(z_j)$, and the generator acts by
\[
 (x,s_j)\longmapsto(A_jx+v_j/m_j,
 e^{-2\pi i/m_j}s_j).
\]
Radial contraction of the disc is equivariant.  The universal cover of
$S_j$ is $\Lambda\otimes\R$, and the group generated by translations
$t_\lambda$ and the affine lift
$x\mapsto A_jx+v_j/m_j$ has the displayed relations; its $m_j$-th power
is $t_{v_j}$.  On the boundary, the lift of one $m_j$-th of the clockwise
base circle is an infinite-order element $\wh g_j$.  Traversing it
$m_j$ times gives the lattice translation $v_j$ together with the full
boundary circle of the covering disc.  Killing that boundary circle on
filling the disc gives exactly the normal relation
$\wh g_j^{m_j}t_{-v_j}$.
\end{proof}

\subsection{Gluing}
Shrink $D_0,D_1,D_2$ so that they are pairwise disjoint and
\[
 B=B^\circ\cup D_0\cup D_1\cup D_2.
\]
Use the biholomorphism $G_0$ of \cref{prop:cusp-identification} and the
biholomorphisms $G_1,G_2$ of \cref{thm:finite-fillings}.

Define
\[
 X=(J\sqcup N_0\sqcup N_1\sqcup N_2)/\sim,
\]
where points in $J|_{D_0^*}$ are identified through $G_0$, and points in
$N_j\setminus S_j$ through $G_j$.

\begin{theorem}[The compact complex threefold]\label{thm:global-X}
The space $X$ is a connected compact Hausdorff complex threefold.  The
local maps glue to a proper surjective holomorphic map
\[
 f:X\longrightarrow B=\PP^1
\]
with connected fibres.  Over $B^\circ$ its fibres are complex two-tori;
$f^{-1}(p_0)=W$ is the reduced normal crossings surface of
\cref{thm:cusp-filling}; and
\[
 f^*(p_1)=3S_1,
 \qquad
 f^*(p_2)=4S_2.
\]
\end{theorem}

\begin{proof}
The three overlap regions in $J$ are disjoint.  Hence every nonsingleton
equivalence class consists of exactly two points, and the quotient maps
from the four pieces are open embeddings.  Their transition maps are the
biholomorphisms $G_0^{\pm1},G_1^{\pm1},G_2^{\pm1}$, so they define a
complex atlas on $X$.  Since the four pieces are second countable and their
images are open, $X$ is second countable.

Points with different images in $B$ are separated by inverse images of
disjoint base neighbourhoods.  Points over the same base point lie in one
of the Hausdorff local models, proving that $X$ is Hausdorff.  Properness
is local on the base and holds for each of the four local maps; therefore
$f$ is proper and $X=f^{-1}(B)$ is compact.  Connectedness and the fibre
descriptions follow from the corresponding properties of the local
models.
\end{proof}

